\begin{trapbox}
\textbf{误区一：沿用椭圆的“大分母决定焦点”。} 双曲线应看正项，而不是分母大小。
\begin{center}\begin{tikzpicture}[font=\small,>=stealth]
\draw[black!30,rounded corners=2pt] (0,0) rectangle (14.4,8.2);
\node[anchor=west,font=\small\bfseries] at (.5,7.72) {看正项确定实轴：双曲线不用比较分母大小};
\begin{scope}[shift={(3.6,4.2)},scale=.4]
\draw[->,black!40] (-6,0)--(6,0) node[right] {$x$};\draw[->,black!40] (0,-6)--(0,6) node[above] {$y$};
\draw[blue,very thick,domain=-4:4,samples=121,variable=\t] plot ({-1*3*sqrt(1+\t*\t/16)},\t);
\draw[blue,very thick,domain=-4:4,samples=121,variable=\t] plot ({1*3*sqrt(1+\t*\t/16)},\t);
\fill (-5,0) circle (2pt) node[below] {$F_1$};
\fill (5,0) circle (2pt) node[below] {$F_2$};
\end{scope}
\node at (3.6,1.3) {$\frac{x^2}{9}-\frac{y^2}{16}=1$};
\begin{scope}[shift={(10.8,4.2)},scale=.4]
\draw[->,black!40] (-6,0)--(6,0) node[right] {$x$};\draw[->,black!40] (0,-6)--(0,6) node[above] {$y$};
\draw[blue,very thick,domain=-4:4,samples=121,variable=\t] plot (\t,{-1*3*sqrt(1+\t*\t/16)});
\draw[blue,very thick,domain=-4:4,samples=121,variable=\t] plot (\t,{1*3*sqrt(1+\t*\t/16)});
\fill (0,-5) circle (2pt) node[right] {$F_1$};
\fill (0,5) circle (2pt) node[right] {$F_2$};
\end{scope}
\node at (10.8,1.3) {$\frac{y^2}{9}-\frac{x^2}{16}=1$};

\node[anchor=west] at (.7,.55) {同尺度；两图均取 $a=3,\ b=4,\ c=5$，焦点位于开口方向};
\end{tikzpicture}\end{center}

两图 $b=4>a=3$，但焦点仍沿正项对应的坐标轴。$a,b$ 不要求谁大，$c$ 才是三者中最大的。

\textbf{误区二：省去绝对值仍表示两支。}
必须写两个等式 $r_1-r_2=\pm2a$，或写一个绝对值等式；只写正号会漏掉左支。

\textbf{误区三：焦点三角形总是非退化的。}
顶点 $P=(\pm a,0)$ 与两焦点共线，距离和此时可等于焦距；普遍使用的应是非严格三角不等式。
需要求非零面积或角度时，再明确 $P$ 不为顶点。
\end{trapbox}
\clearpage
